\scanpage{43}\mspage{42}
\uline{caractérisée} par les conditions suivantes :
% slanted
\margin{C'est faux, ça implique (i) (ii) (iii) (iv), mais ces conditions ne suffisent nullement, voir plus bas ! \unsure{p.|v.}~48}
\margin{i.e $\overline{{\unsure{\mathcal{T}|\tau}}'}\in\underline{\underline{\overline{\unsure{\mathscr{T}|\mathcal{T}}}}}$ $\left\{\vphantom{\rule{0pt}{4ex}}\right.$}

(i) $\overline{\mathcal{T}'}$ est plus fine que $\widetilde{\mathcal{T}}$, et induit sur $J'$ $\unsure{\widetilde{\mathcal{T}}|\overline{\mathcal{T}}}\mid J'$.

(ii) $\overline{\mathcal{T}'}\mid\overset{\deleted{\unsure{\sim|\mathord{-}}}}{{\unsure{U|u|n}}'}=I'\setminus J'$ coïncide avec $\widetilde{\mathcal{T}}\mid\overset{\deleted{\sim}}{U'}$

(iii) $\varphi\mid U'$ induit \deleted{\unsure{une bijection}} \inserted{\unsure{iso. préord.}} $U'\longrightarrow U=S\setminus T$
% slanted
\margin{\deleted{\unsure{Autant dire|Autrement dit}}}

(\unsure{iv|iiv}) La topologie de $I$ est quotient \unsure{final|finale|déduit} pour les applications $I'\unsure{\mathbin{\text{---}}|\longrightarrow}I$, $J\mathbin{\text{---}}I$ ($I'$ munie \deleted{\unsure{celle de} $I$, \unsure{conjointement}} \inserted{de $\overline{\mathcal{T}'}$, $J$ \deleted{\unsure{de}} muni de $\overline{\mathcal{T}}\mid J$}) (\unsure{Ce} qui exclut notamment la solution triviale pour (i) (ii), obtenue en prenant la topologie somme sur \deleted{\unsure{$J'$}} $I'=J'\underset{\text{disj.}}{\cup}\overset{\deleted{\sim}}{U'}$, sauf si $I$ lui\unsure{$=$} est une somme top. $J\sqcup\underset{\substack{\Vert\\I\setminus J}}{U}$).
% slanted
\margin{C'est cette condition que j'avais d'abord \unsure{oubliée|oublié}, à la page 31 !}
% slanted
\margin{\deleted{\unsure{Autant}} dire que $\mathcal{T}'$ sur $S$ est caractérisée par les propriétés analogues (i) (ii) (iii) (iv), la dernière \unsure{signifiant|signifie} que la top $\text{\textcircled{$\mathcal{T}$}}$ sur $S$ est finale pour $\unsure{\mathcal{T}''|\mathcal{T}'}$ sur $S$ et $\mathcal{T}$ sur $T$ (\unsure{id|et} $S\unsure{\rightrightarrows|\to}S$ $T\unsure{\hookrightarrow|\to}S$)}
% slanted
\margin{en fait, c'est indépendant \unsure{aussi|même} de l'hyp. \unsure{ultra-topologique}}

Démonstration

\deleted{\unsure{a) et b).}} \inserted{(a b c)} Il est clair que les ens. \unsure{épinglés} dans c) sont ceux qui satisfont $1^\circ)$, $2^\circ)$ $3^\circ)$ dans b), et que la représentation dans c) est unique. Il reste à voir, pour prouver (a,b,c), que les ens. forment une topologie, \deleted{et} qu'elle \deleted{\unsure{a les propriétés}} \inserted{est} dans $\underline{\mathcal{T}}$, et est moins fine des top $\in\underline{\mathcal{T}}$.

1) Ce sont les fermés d'une topologie.

Stables par intersections quelconques : clair

Stables par réunions finies : idem.

C'est \unsure{$=$} stable par réunions \unsure{quelconques} s'il en est ainsi \unsure{pour|dans} $(T,\mathcal{T}'_0)$, $(\deleted{\unsure{D|I}}\unsure{S|I},\mathcal{T})$, d'où la première assertion de \unsure{d)|a)}.

